% !TeX root = ../../Book2.tex
\OptionalSection{Besov 与 Triebel–Lizorkin 空间}{b2:sec:1905}

同样一组频率块，可以记录两种不同的信息。先量每一块在整个空间中的大小，再比较不同尺度，得到一个数列范数；先在每个位置比较各尺度，再对位置积分，得到另一种范数。\subsecref{b2:subsec:190404}的平方空间已经采用了后一种次序。本节允许尺度指数变化，并通过差分、乘积和平坦边界的迹，说明这个指数为什么不能省略。一般空间可以选读；\chapref{b2:ch:22}的普通乘子正文只需要上一节已经建立的平方情形。

\subsection{二进定义}
\label{b2:subsec:190501}

先把分割固定下来，使后面的每个范数都有确定含义。低频块的作用也必须保留：所有齐次环带都消去多项式，而非齐次分解还能记录零频附近的信息。

\Needspace{6\baselineskip}
\begin{definition}\label{b2:def:19-besov-triebel-lizorkin}
设整数 \(d\geq1\)，取实径向函数 \(\chi\in C_c^\infty(\mathbb R^d)\)，满足 \(0\leq\chi\leq1\)，在 \(|\xi|\leq1\) 为一、\(|\xi|\geq2\) 为零，并沿半径单调不增。令 \(\varphi=\chi-\chi(2\,\cdot)\)。对复值温和分布 \(f\in\mathcal S'(\mathbb R^d)\)，定义
\[
 P_0f=\mathcal F^{-1}(\chi\cdot\widehat f),
 \quad P_jf=\mathcal F^{-1}\bigl(\varphi(2^{-j}\,\cdot)\cdot\widehat f\bigr)
 \quad(j\geq1).
\]
设 \(s\in\mathbb R\)、\(1\leq p,q\leq\infty\)。满足
\[
 \|f\|_{B^s_{p,q}}
 \coloneqq\left\|\bigl\{2^{js}\|P_jf\|_p\bigr\}_{j\geq0}\right\|_{\ell^q}<\infty
\]
的温和分布组成\term[besov-kongjian]{Besov 空间}[Besov space]，记为 \(B^s_{p,q}(\mathbb R^d)\)。当 \(1\leq p<\infty\)、\(1\leq q\leq\infty\) 时，满足
\[
 \|f\|_{F^s_{p,q}}
 \coloneqq\left\|\left\|\bigl\{2^{js}P_jf(x)\bigr\}_{j\geq0}\right\|_{\ell^q}\right\|_{L^p_x}<\infty
\]
的温和分布记为 \(F^s_{p,q}(\mathbb R^d)\)，称为一般的 Triebel--Lizorkin 空间。序列指数 \(q=\infty\) 时，\(\ell^q\) 范数取上确界；Besov 定义中的 \(p=\infty\) 取空间本性上确界。
\end{definition}
\symidx{\(B^s_{p,q}\)：非齐次 Besov 空间}
\symidx{\(F^s_{p,q}\)：非齐次 Triebel--Lizorkin 空间}
这里的块正是\cref{b2:def:19-inhomogeneous-square}中的块。每个块都是光滑函数，但未必有有限的全空间范数；定义的有限性不是自动得到的。由于
\[
 P_0f+\sum_{j=1}^NP_jf
 =\mathcal F^{-1}\bigl(\chi(2^{-N}\,\cdot)\cdot\widehat f\bigr)
 \longrightarrow f\quad\text{于 }\mathcal S',
\]
这些范数为零时所有块为零，继而 \(f=0\)。三角不等式分别来自 \(L^p\) 与 \(\ell^q\) 的三角不等式，所以在本节的指数范围内它们确实是范数。

\ProseParagraphBreak

三个参数分别回答三个问题：\(s\) 规定高频块要乘多大的尺度权；\(p\) 规定怎样量空间大小；\(q\) 规定允许多少个尺度同时显著。只写光滑阶数而不写后两个指数，通常不能判断函数是否属于该空间。

\subsection{Sobolev 空间的重新解释}
\label{b2:subsec:190502}

\Cref{b2:prop:19-inhomogeneous-identification}已证明
\[
 F^0_{p,2}(\mathbb R^d)=L^p(\mathbb R^d)
 \quad(1<p<\infty),
 \quad F^s_{2,2}(\mathbb R^d)=H^s(\mathbb R^d)
 \quad(s\in\mathbb R),
\]
等号指空间相同、范数等价。前一识别包括反向的函数代表问题：若温和分布的非齐次平方函数属于 \(L^p\)，合成估计使其块和在 \(L^p\) 中收敛，分布重构再将极限识别为原分布。不能只对已知属于 \(L^p\) 的函数证明不等式，就宣布对所有温和分布完成了空间识别。

\ProseParagraphBreak

在 \(p=q=2\) 时，Tonelli 允许交换位置积分与尺度求和，直接给出 \(B^s_{2,2}=F^s_{2,2}=H^s\)。一般 \(p=q<\infty\) 也有定义范数完全相等的恒等式
\[
 \sum_{j\geq0}2^{jsp}\|P_jf\|_p^p
 =\int_{\mathbb R^d}\sum_{j\geq0}2^{jsp}|P_jf(x)|^p\,\mathrm dx.
\]
这一次是两个积分次序的交换；\(H^s\) 与频块范数的比较还涉及频率权重，所以通常只有等价而不是数值相等。

\ProseParagraphBreak

另一个常用描述从符号 \((1+4\pi^2|\xi|^2)^{s/2}\) 定义 \(J^s\)，再要求 \(J^sf\in L^p\)。该符号光滑、各阶导数至多作多项式增长，因而可作用于温和分布。对于 \(1<p<\infty\)，把这个条件与 \(F^s_{p,2}\) 识别，需要统一控制环带上的归一化符号及其逆；这一乘子论证见\cref{b2:prop:22-bessel-square}。本章的平方可积情形已由 Plancherel 独立完成。齐次势 \(I_\alpha\) 在零频奇异，而 \(J^{-\alpha}\) 的符号在那里有限；这种低频差别正是全空间同指数估计不同的原因，不能把二者混作同一个算子。

\subsection{Besov 空间}
\label{b2:subsec:190503}

先说明所选分割如何影响 Besov 范数。若 \((\widetilde P_k)\) 来自另一套满足定义条件的截断，则高频新块只与固定个数的旧块相交，低频处也只有固定个数的例外。分布重构因此给出有限和
\[
 \widetilde P_kf=\sum_{j\sim k}\widetilde P_kP_jf.
\]
这里 \(j\sim k\) 表示 \(|j-k|\leq L\)，可增大固定整数 \(L\) 以覆盖低频例外。新块卷积核的 \(L^1\) 范数一致有界，且在这些指标上 \(2^{ks}\) 与 \(2^{js}\) 可比，故
\[
 2^{ks}\|\widetilde P_kf\|_p
 \leq C_s\sum_{j\sim k}2^{js}\|P_jf\|_p.
\]
对序列取 \(\ell^q\) 范数，有限次平移不改变该范数，便得到一个方向的界；交换新旧分割即得反向界。这个证明覆盖 \(1\leq p,q\leq\infty\)，常数依赖两套固定截断。

\ProseParagraphBreak

为了具体看见第二指数，取 \(s>0\)，在彼此相隔固定个数二进尺度的频率小球内选 \(f_n\in\mathcal S\)，使 \(\|f_n\|_2=1\)。记相应尺度为 \(j_n=Ln\)，其中固定 \(L\) 足够大。令
\[
 f=\sum_{n\geq1}a_n2^{-j_ns}f_n,
\]
并先设 \((a_n)\) 有界。不同球互不相交，正 \(s\) 使级数在 \(L^2\) 中收敛。每个球只与固定个数的频块相交，不同球对应的指标组互不交叠；在每组中，分割和为一，至少有一个块记录了固定比例的平方大小。Plancherel 和有限个数项的范数等价因此给出
\[
 \|f\|_{B^s_{2,q}}\asymp\|(a_n)\|_{\ell^q}.
\]
于是 \(a_n=1/n\) 给出属于 \(B^s_{2,2}\) 而不属于 \(B^s_{2,1}\) 的函数；\(a_n=1\) 给出属于 \(B^s_{2,\infty}\) 而不属于 \(H^s\) 的函数。频率位置相同的尺度权并不能替代尺度之间的可和性。

\ProseParagraphBreak

当光滑阶数小于一时，还可以不用 Fourier 变换来辨认一种 Besov 范数。频率块在步长远小于其波长时变化较小，在步长较大时只能用两倍振幅控制；下面的两段几何级数正来自这一区别。

\Needspace{6\baselineskip}
\begin{proposition}\label{b2:prop:19-besov-difference}
设整数 \(d\geq1\)、\(0<s<1\)、\(1\leq p<\infty\)，采用\cref{b2:def:19-besov-triebel-lizorkin}的非齐次块。对每个 \(f\in L^p(\mathbb R^d;\mathbb C)\)，有
\[
 \|f\|_{B^s_{p,p}}\asymp
 \|f\|_p+
 \left(\int_{\mathbb R^d}
 \frac{\|f(\,\cdot+h)-f\|_p^p}{|h|^{d+sp}}\,\mathrm dh\right)^{1/p}.
\]
允许两边同时为无穷；比较常数只依赖 \(d,s,p\) 与固定截断。因此该 Besov 空间与按右侧差分范数定义的分数 Sobolev 空间相同。
\end{proposition}
\begin{proof}
先设左边有限，记 \(u_j=P_jf\)、\(a_j=2^{js}\|u_j\|_p\)。由于 \(s>0\)，Hölder 的序列形式（\(p=1\) 时直接求和）给出
\[
 \sum_{j\geq0}\|u_j\|_p
 =\sum_{j\geq0}2^{-js}a_j
 \leq C_{s,p}\|(a_j)\|_{\ell^p}.
\]
所以块和在 \(L^p\) 中绝对收敛；分布重构说明其和就是 \(f\)，并且已经得到所需的 \(\|f\|_p\) 界。

沿线段积分及 Bernstein 估计给出 \(\|u_j(\cdot+h)-u_j\|_p\leq C2^j|h|\cdot\|u_j\|_p\)；平移不变性另给出两倍函数范数的界。因此
\[
 \|u_j(\,\cdot+h)-u_j\|_p
 \leq C\min\{1,2^j|h|\}\cdot\|u_j\|_p.
\]
在 \(2^{-k-1}<|h|\leq2^{-k}\)、\(k\geq0\) 上，对绝对重构求差并相加，得到
\[
 2^{ks}\|f(\,\cdot+h)-f\|_p
 \leq C\sum_{j\leq k}2^{-(k-j)(1-s)}a_j
       +C\sum_{j>k}2^{-(j-k)s}a_j.
\]
右边是两段可和几何序列与 \((a_j)\) 的卷积，离散 Young 不等式控制其 \(\ell^p_k\) 范数。每个步长环带上的 \(\int|h|^{-d}\,\mathrm dh\) 相同，且 \(|h|^{-sp}\) 与 \(2^{ksp}\) 可比，所以这个序列界正给出 \(|h|\leq1\) 的差分积分。\(|h|>1\) 的部分由 \(2\|f\|_p\) 以及 \(\int_{|h|>1}|h|^{-d-sp}\,\mathrm dh<\infty\) 控制。

反向设右边有限。对 \(j\geq1\)，令 \(k_j(h)=2^{jd}\check\varphi(2^jh)\)。其积分为零，因此在 \(L^p\) 中有
\[
 P_jf=\int k_j(h)\cdot\bigl(f(\,\cdot-h)-f\bigr)\,\mathrm dh.
\]
Minkowski 后对测度 \(|k_j(h)|\,\mathrm dh\) 使用 Hölder，得到
\[
 \|P_jf\|_p^p
 \leq\|k_j\|_1^{p-1}
       \int|k_j(h)|\cdot\|f(\,\cdot-h)-f\|_p^p\,\mathrm dh.
\]
\(p=1\) 时同式直接成立。Schwartz 衰减又给出
\[
 \sum_{j\geq1}2^{jsp}|k_j(h)|\leq C|h|^{-d-sp}\quad(h\ne0).
\]
具体地，写 \(r=|h|\)，每项至多为 \(Cr^{-d-sp}(2^jr)^{d+sp}(1+2^jr)^{-N}\)。在 \(2^jr\leq1\) 与 \(>1\) 两侧分别求几何级数，取 \(N>d+sp\) 即可。乘以差分范数后用 Tonelli，便控制全部高频项；低频块用 \(\|P_0f\|_p\leq\|\check\chi\|_1\cdot\|f\|_p\)。两向估计由 \(L^p\) 差分与各频带的范数联系起来。
\end{proof}
\(s<1\) 保证小步长对低频块的几何衰减可和，\(s>0\) 保证高频重构与大步长积分收敛。若 \(s\geq1\)，一般必须换用更高阶差分；对非零光滑梯度，\(s=1\) 的一阶差分积分已经出现 \(\mathrm dr/r\) 的发散。

\ProseParagraphBreak

这个判别可以在一个跳跃函数上精算。取单位立方体 \(Q=(0,1)^d\)。当每个 \(|h_i|<1\) 时，交集体积为 \(\prod_i(1-|h_i|)\)，因而对任意有限 \(p\geq1\)，
\[
 \|\mathbf1_Q(\,\cdot+h)-\mathbf1_Q\|_p^p
 =|Q\mathbin{\triangle}(Q-h)|
 =2\left(1-\prod_{i=1}^d(1-|h_i|)\right).
\]
若 \(|h|<1/2\)，括号介于 \(\max_i|h_i|\) 与 \(\sum_i|h_i|\) 之间，故与 \(|h|\) 可比。小步长积分于是与 \(\int_0^{1/2}r^{-sp}\,\mathrm dr\) 同时收敛。得到精确条件
\[
 \mathbf1_Q\in B^s_{p,p}\quad\Longleftrightarrow\quad sp<1
 \quad(0<s<1,\ 1\leq p<\infty).
\]
区域内外的函数都很简单，限制光滑阶数的是跨越边界时保留的跳跃。

\ProseParagraphBreak

绝对尺度求和还控制函数乘积。由 Bernstein，
\[
 \sum_{j\geq0}\|P_jf\|_\infty
 \leq C\sum_{j\geq0}2^{jd/p}\|P_jf\|_p
 =C\|f\|_{B^{d/p}_{p,1}}.
\]
所以 \(B^{d/p}_{p,1}\) 有一个有界连续代表：每块连续，上确界范数绝对可和，故块和一致收敛；分布重构识别这个代表。这也说明为什么临界处的序列指数一有特殊作用。

\Needspace{6\baselineskip}
\begin{proposition}\label{b2:prop:19-besov-algebra}
设整数 \(d\geq1\)、\(1\leq p<\infty\)、\(s\geq d/p\)，采用\cref{b2:def:19-besov-triebel-lizorkin}的非齐次 Besov 空间。则它的元素具有有界连续代表，其逐点乘积满足
\[
 \|fg\|_{B^s_{p,1}}
 \leq C_{d,s,p,\chi}\|f\|_{B^s_{p,1}}\cdot\|g\|_{B^s_{p,1}}.
\]
特别地，\(B^s_{p,1}(\mathbb R^d)\) 在乘法下封闭。
\end{proposition}
\begin{proof}
令 \(f_k=P_kf\)、\(g_l=P_lg\)。因为 \(s\geq d/p>0\)，两组块都在 \(L^p\) 和 \(L^\infty\) 中绝对求和，前述嵌入给出有界连续代表。双级数 \(\sum_{k,l}f_kg_l\) 在 \(L^p\) 中绝对收敛，并等于 \(fg\)。

每项乘积的 Fourier 支集包含于两频率球之和，故在半径 \(C2^{\max\{k,l\}}\) 的球内。这一支集规则可以对紧支集频率分布作卷积来理解；也可先用相同频率截断将 \(L^p\) 光滑逼近保持在略大的固定频带内，再把 Schwartz 函数的卷积公式传到分布极限。因而存在固定整数 \(L\)，使 \(P_j(f_kg_l)=0\) 当 \(j>\max\{k,l\}+L\)。

若 \(k\geq l\)，由块核的统一 \(L^1\) 界及 \(s>0\)，
\[
 \begin{aligned}
 \sum_{j\geq0}2^{js}\|P_j(f_kg_l)\|_p
 &\leq C\sum_{0\leq j\leq k+L}2^{js}\|f_k\|_p\cdot\|g_l\|_\infty\\
 &\leq C_s2^{ks}\|f_k\|_p\cdot\|g_l\|_\infty.
 \end{aligned}
\]
\(l>k\) 时交换两因子。先对有限双和估计，再用 \(L^p\) 收敛及 Fatou，得到
\[
 \begin{aligned}
 \|fg\|_{B^s_{p,1}}
 \leq C_s&\left(\sum_k2^{ks}\|f_k\|_p\right)
          \left(\sum_l\|g_l\|_\infty\right)\\
 +C_s&\left(\sum_k\|f_k\|_\infty\right)
          \left(\sum_l2^{ls}\|g_l\|_p\right).
 \end{aligned}
\]
两组上确界范数和均由对应 Besov 范数控制，即得结论。证明只用了 \(s\geq d/p\) 与 \(s>0\)，所以包括 \(s=d/p\) 的临界情形。
\end{proof}
若两个相乘的块频率相差很大，输出仍以较高频为主；若二者频率接近，却可能相互抵消而产生低频。上面求和的是所有 \(j\leq\max\{k,l\}+L\)，已经把这些低频输出包括在内，不能只保留高低频相互作用。这项乘积结论依赖已得到的函数代表，不能推及任意两个温和分布。

\subsection{Triebel--Lizorkin 空间}
\label{b2:subsec:190504}

这类范数先在位置 \(x\) 处比较不同频率，再作空间积分。若许多块集中在同一小区域，点态序列范数会同时记录它们；若这些块分布在彼此远离的区域，同一位置只见到少数项。先把每块的全空间范数算成一个数，会丢掉这种位置关系。两种次序可以比较，却通常不能交换。

\Needspace{6\baselineskip}
\begin{proposition}\label{b2:prop:19-besov-triebel-embeddings}
设整数 \(d\geq1\)、\(s\in\mathbb R\)、\(1\leq p<\infty\)、\(1\leq q\leq\infty\)，所有空间使用\cref{b2:def:19-besov-triebel-lizorkin}中的同一分割。则有连续嵌入
\[
 B^s_{p,\min\{p,q\}}\subset F^s_{p,q}
 \subset B^s_{p,\max\{p,q\}}.
\]
此外，对 \(1\leq p\leq r\leq\infty\)、\(1\leq q\leq\infty\)、\(s\in\mathbb R\)，有
\[
 B^s_{p,q}\subset B^{s-d(1/p-1/r)}_{r,q},
\]
其中 \(1/\infty=0\)。
\end{proposition}
\begin{proof}
先对有限个块令 \(a_j(x)=2^{js}|P_jf(x)|\)。若 \(q\leq p\)，Minkowski 在 \(L^{p/q}\) 中给出
\[
 \left\|\sum_j a_j^q\right\|_{p/q}
 \leq\sum_j\|a_j^q\|_{p/q},
\]
即 \(\|f\|_{F^s_{p,q}}\leq\|f\|_{B^s_{p,q}}\)。逐点的 \(\ell^p\leq\ell^q\) 又给出 \(\|f\|_{B^s_{p,p}}\leq\|f\|_{F^s_{p,q}}\)。若 \(p\leq q<\infty\)，Minkowski 在序列空间 \(\ell^{q/p}\) 中给出
\[
 \left(\sum_j\left[\int a_j^p\right]^{q/p}\right)^{p/q}
 \leq\int\left(\sum_j a_j^q\right)^{p/q},
\]
从而 \(\|f\|_{B^s_{p,q}}\leq\|f\|_{F^s_{p,q}}\)；逐点的 \(\ell^q\leq\ell^p\) 给出另一边。\(q=\infty\) 时直接使用
\[
 \sup_j\|a_j\|_p\leq\|\sup_j a_j\|_p
 \leq\left(\sum_j\|a_j\|_p^p\right)^{1/p}.
\]
有限块数趋于无穷，单调收敛证明第一组嵌入。

第二组嵌入对每个块使用 Bernstein：
\[
 2^{j[s-d(1/p-1/r)]}\|P_jf\|_r
 \leq C2^{js}\|P_jf\|_p.
\]
低频块尺度为一，满足同样估计。取 \(\ell^q\) 范数即得结论，包括两个无穷端点。
\end{proof}
这些都是范数不等式；只有另有严格性例子时，才可把包含号改成真包含。\(p=q\) 时第一组嵌入就是前面由 Tonelli 得到的相等，因而不可能一概宣称严格。

\ProseParagraphBreak

分割独立性也要区分指数。对 \(q=2\)、\(1<p<\infty\)，写
\(\widetilde P_kf=\sum_{|j-k|\leq L}\widetilde P_kP_jf\)，并将有限个相邻偏移分别处理。每个固定偏移对应一族 \(\widetilde P_k\)：其缩放核有统一的标准核常数，平方范数一致有界，所以\cref{b2:prop:18-hilbert-valued-cz}控制 \(L^p(\ell^2)\) 范数。权重 \(2^{ks}\) 与 \(2^{js}\) 在相邻指标上可比，得到新旧 \(F^s_{p,2}\) 范数等价。\(p=1\) 或一般 \(q\) 不在这一 Hilbert 值强型论证的范围内。这里的定义仍使用固定分割，前面的差分、嵌入和例子也均已在声明的分割下证明；一般 Triebel--Lizorkin 分割独立性的进一步理论可参看 Hans Triebel 的《Theory of Function Spaces》\cite{Triebel1983Theory}，不能用标量 Young 不等式代替其中的向量估计。

\subsection{Besov 空间的平面迹与延拓}
\label{b2:subsec:190505}

最后考虑在 \(x_d=0\) 上取边界值。一个频率为 \(2^j\) 的块，在法向上可变化的长度约为 \(2^{-j}\)。把这一层的 \(d\) 维积分变成边界的 \((d-1)\) 维积分，函数范数会出现 \(2^{j/p}\)，因此 \(s\) 阶的自然边界阶数为 \(s-1/p\)。下面把这项尺度计算落实为一个有界迹与右逆，范围限定为平坦边界、序列指数一。

\Needspace{6\baselineskip}
\begin{proposition}[平坦边界的 Besov 迹与延拓]
\label{b2:prop:19-flat-besov-trace}
设整数 \(d\geq2\)、\(1\leq p<\infty\)、\(s>1/p\)。在 \(\mathbb R^d\) 与 \(\mathbb R^{d-1}\) 上各固定\cref{b2:def:19-besov-triebel-lizorkin}的一套非齐次分割。Schwartz 函数上的限制 \(u\mapsto u(\,\cdot,0)\) 有唯一有界线性延拓
\[
 \operatorname{Tr}:B^s_{p,1}(\mathbb R^d)
 \longrightarrow B^{s-1/p}_{p,1}(\mathbb R^{d-1}),
\]
并存在有界线性算子 \(E\) 从目标空间映回 \(B^s_{p,1}(\mathbb R^d)\)，满足 \(\operatorname{Tr}Eg=g\)。各常数只依赖固定参数和截断。
\end{proposition}
\begin{proof}
令 \(u_j=P_ju\)。其法向频率不超过 \(C2^j\)。取一维 Schwartz 函数 \(k\)，使 \(\widehat k=1\) 于 \([-C,C]\)，则由法向频率截断，
\[
 u_j(x',0)=\int_{\mathbb R}2^jk(-2^jt)\cdot u_j(x',t)\,\mathrm dt.
\]
对每个 \(x'\) 使用 Hölder，\(p'=p/(p-1)\) 给出核的 \(L^{p'}\) 范数 \(2^{j/p}\|k\|_{p'}\)；\(p=1\) 时使用 \(p'=\infty\)。再积分切向变量，得到
\[
 \|u_j(\,\cdot,0)\|_{L^p(\mathbb R^{d-1})}
 \leq C2^{j/p}\|u_j\|_{L^p(\mathbb R^d)}.
\]
卷积表示给出所用的法向连续代表；频率块本身光滑，因此这就是其普通限制。边界频率仍在半径 \(C2^j\) 的球内。记边界块为 \(P'_k\)，其核的 \(L^1\) 范数一致有界，而且 \(P'_k[u_j(\cdot,0)]=0\) 当 \(k>j+L\)。于是
\[
 \begin{aligned}
 \|u_j(\,\cdot,0)\|_{B^{s-1/p}_{p,1}}
 &\leq C\sum_{0\leq k\leq j+L}2^{k(s-1/p)}
              \cdot2^{j/p}\|u_j\|_p\\
 &\leq C_s2^{js}\|u_j\|_p.
 \end{aligned}
\]
最后的几何级数使用了 \(s-1/p>0\)。对 \(j\) 求和，边界块的限制之和在 \(L^p\) 中绝对收敛，记为 \(\operatorname{Tr}u\)；逐个边界块取极限并用 Fatou，得到目标 Besov 界。

补足延拓的唯一性。因为序列指数为一，有限频率和在 \(B^s_{p,1}\) 中稠密：对剩余块使用有限重叠及统一卷积界，其范数受原尺度序列的尾和控制。若一个有限频率和 \(v\in L^p\) 的频率支集在固定球内，选光滑紧频率乘子 \(\eta\)，在该球邻域内为一。取 \(h_n\in C_c^\infty\) 在 \(L^p\) 中逼近 \(v\)，则 \(\eta(D)h_n\in\mathcal S\)，且
\[
 \|\eta(D)h_n-v\|_p
 \leq\|\check\eta\|_1\cdot\|h_n-v\|_p\longrightarrow0.
\]
误差有共同有界频带，故只含有限个非齐次块，\(L^p\) 收敛也给出 \(B^s_{p,1}\) 收敛。由此证明 Schwartz 类稠密。对 Schwartz 函数，非齐次截断在 Schwartz 拓扑中趋于原函数，边界限制与所定义的块和相同，所以所得连续延拓唯一。

为构造右逆，将 \(g\) 分解为边界块 \(g_j=P'_jg\)。选一维 Schwartz 函数 \(\rho\)，使 \(\rho(0)=1\) 且 \(\widehat\rho\) 紧支撑，令
\[
 v_j(x',t)=g_j(x')\cdot\rho(2^jt),
 \quad Eg=\sum_{j\geq0}v_j.
\]
每项 Fourier 支集在半径 \(C2^j\) 的球内，且
\[
 \|v_j\|_{L^p(\mathbb R^d)}
 =2^{-j/p}\|g_j\|_{L^p(\mathbb R^{d-1})}\cdot\|\rho\|_{L^p(\mathbb R)}.
\]
对体内输出块再求和，由 \(s>0\) 得
\[
 \|v_j\|_{B^s_{p,1}(\mathbb R^d)}
 \leq C_s2^{js}\|v_j\|_p
 \leq C2^{j(s-1/p)}\|g_j\|_p.
\]
右侧可和，所以 \(\sum v_j\) 在 \(L^p\) 中收敛，并由块估计与 Fatou 得到 \(Eg\in B^s_{p,1}\) 及相应范数界；同一尾和估计也证明级数在 Besov 范数中收敛。每个 \(v_j\) 只有有限多个体内频率块，其块和就是本身，故由定义 \(\operatorname{Tr}v_j=v_j(\cdot,0)=g_j\)。连续性将有限和等式传给极限，得到 \(\operatorname{Tr}Eg=\sum_jg_j=g\)。
\end{proof}
临界 \(s=1/p\) 时，边界输出尺度的求和不再有几何衰减，不能由本证明得到同样的 Besov 目标；但是只求边界 \(L^p\) 范数时，绝对求和还能保留一个临界版本，见\cref{b2:ex:19-new-3}。曲边界则另需坐标变换的有界性，不能把平面上的公式直接搬到任意区域。迹的尺度损失、序列指数与边界几何，各自承担不同的条件。
